\chapter{Literature Review and Existing Systems}
\label{ch:literature-review}

\begin{center}
\textit{This chapter reviews existing systems and technologies relevant to
the platform --- learning management systems, optical character recognition,
large language models in education, multi-tenancy, authorization models, and
asynchronous processing patterns --- and identifies the gap that the project
addresses.}
\end{center}

\vspace*{\fill}
\clearpage
\section{Existing Educational Platforms}

Learning management systems (LMS) such as Moodle~\cite{moodle}, Google
Classroom~\cite{gclassroom}, and Canvas~\cite{canvas} provide course delivery,
assignments, and basic content organization to whole institutions. They are
strong at the operational layer --- announcing courses, distributing
assignments, and tracking participants --- but expose limited support for the
specific preparation-and-assessment workflow this project targets: turning
uploaded textbook chapters and previous-year papers into structured content,
a structured question bank aligned to an exam blueprint, and scheduled
server-evaluated assessments. Such platforms do not typically perform OCR of
arbitrary educational documents, nor do they generate study content and
examination assets with an AI pipeline and provenance for every produced
artifact.

Commercial and open-source tools exist for individual pieces of the workflow.
OCR engines such as Tesseract and PaddleOCR extract text from scanned pages;
document libraries such as PyMuPDF give programmatic access to PDF text and
page rasterisation; AI content tools produce text on demand; quiz systems
evaluate objective questions. The gap addressed here is a single
tenant-isolated system that \emph{integrates} these pieces with an explicit
authorization model, so that a college gets one coherent asset pipeline from
``upload'' to ``evaluate''.

\section{Optical Character Recognition}

Tesseract is a long-established open-source OCR engine with excellent
support for structured Latin-script documents~\cite{tesseract}. For
handwritten and complex scripts, deep-learning OCR models are more robust.
The PaddleOCR line of models (PP-OCR)~\cite{ppocr} provides lightweight,
practical, multilingual OCR pipelines from the PaddlePaddle ecosystem and is
widely used for real-world degradation (noise, watermarks, low contrast).
Document payloads almost always live inside PDF files, which may contain an
embedded digital text layer (born-digital PDFs) or be pure scans. PyMuPDF
gives safe, dependency-light access to both: it extracts embedded text per
page and rasterises pages to images for OCR fallback~\cite{pymupdf}.

The design decision of this project is \emph{embedded-text-first, OCR-fallback
per page}: a page whose embedded text has a usable length (at least five
characters) is accepted directly, and only otherwise is it rasterised and
passed to PaddleOCR. This reduces OCR load, preserves born-digital fidelity,
and reserves compute for genuinely scanned material. OCR output is then
normalised by a deterministic character pass, and the resulting text is stored
on the material row so it can feed AI generation and text-based reuse without
re-running vision models.

\section{Large Language Models in Education}

Transformer-based language models~\cite{vaswani2017} have made
instruction-following generation practical for educational content. A
retrieval-augmented pattern --- supply selected source material as context,
then generate the desired artifact (notes, flashcards, summaries, questions)
--- is the dominant integration shape for factual document-derived content
because it grounds the output in the source document and makes the 
\emph{source of truth} explicit~\cite{lewis2020rag}. This project follows that
pattern: the AI worker assembles a bounded context budget from READY source
materials, calls the model through an internal gateway, validates each
response against strict schemas for the requested artifact, and stores
generated items with provenance linking them to the exact source revisions
used. Responses that fail schema validation or transient transport are
retried with backoff up to a maximum number of attempts; terminal errors
settle the associated job as failed rather than leaving it in a running state.

\section{Multi-tenant Architecture}

Multi-tenancy spans a spectrum from isolated deployments (one stack per
tenant) through shared-database models (each tenant's rows prefixed or
separated) with varied data-security trade-offs. The implementation chooses a
shared-schema, tenant-scoped model: one PostgreSQL database whose tables carry
an \texttt{institute\_id}, and every query that touches user data is filtered
by the resolved active institute. Tenant resolution comes from an explicit
header chosen by the client on each request, validated by a guard that checks
a live, active membership. This yields low operational cost and strong
consistency, and it is defensible because the enforcement is at the query
layer rather than the presentation layer.

\section{Authorization Models}

Authorization is classically organised as role-based access control (RBAC),
in which roles bundle rights, or attribute- and policy-based models that
evaluate rules on each request. Pure RBAC is simple but coarse: a ``teacher''
role grants the same rights over every class the teacher teaches, which leaks
access across cohorts. Permission-based models refine RBAC by evaluating the
\emph{permission key} (e.g. \texttt{questions.update}) and, additionally, by
restricting \emph{where} that permission may operate (the academic scope).
Libraries such as CASL provide attribute-based policy evaluation for such
models~\cite{casl}, but this project implements its own compact engine ---- a
single typed permission catalogue with a fixed \texttt{resource.action}
vocabulary, default-deny, a \texttt{manage}-implies action rule, and no
wildcards ---- deliberately avoiding a dependency for a decision that is at
its core a few deterministic predicates. Roles become permission bundles: the
built-in \texttt{INSTITUTE\_ADMIN}, \texttt{TEACHER}, and \texttt{STUDENT}
roles, plus an institution-level \texttt{SUPER\_ADMIN} platform plane that is
kept disjoint from institute membership.

\section{Asynchronous Processing Patterns}

Long-running operation must not block an API request. The classic pattern is a
message broker with durable queues and worker consumers: the API writes a
job row, publishes a message, returns \texttt{202 Accepted} with the job id,
and a worker consumes the message, does the work, and settles the job. This
project uses RabbitMQ~\cite{rabbitmq} with two queues --- one for material
processing and one for AI generation with bounded concurrency. A variant for
the OCR pipeline is a \emph{coordinator-owned} design: the API itself owns the
OCR chunk registry, leases chunks to registered OCR workers by claim, and
reclaims idle or expired leases, because OCR compute nodes are expected to be
flaky, heterogeneous, and occasionally offline. This combination (queue for
message-based work, coordinator for lease-based work) preserves the
API-never-blocks guarantee while tolerating worker loss.

\section{Security Baseline}

Web authentication on this platform follows the OWASP cookie-session
recommendations~\cite{owaspcsfr, owaspsession}: short-lived httpOnly access
cookies, refresh tokens kept server-side as hashes, double-submit CSRF
protection on all state-changing requests, origin checks on credential
submission, and throttling. JWTs~\cite{rfc7519} are used only as
transport-format access tokens and carry no permissions; authorization is
re-evaluated from the current database state per request. Password hashing
uses bcrypt~\cite{bcrypt}, session hashes use SHA-256 per
FIPS~PUB~180-4~\cite{fips180}, and all internal service traffic is carried
over the private Docker network with shared or per-worker keys.

\section{Identified Gap}

Existing systems separately provide LMS features, OCR, AI generation, quiz
evaluation, or multi-tenancy. None of them combine, in one tenant-isolated
product, the end-to-end lifecycle of \emph{OCR-of-materials ---\,
AI-grounding ---\, syllabus structuring ---\, question-bank and paper-pattern
authoring ---\, scheduled server-evaluated assessments ---\, practice ---\,
print-ready export}, under an authorization model that restricts both the
\emph{actions} and the \emph{academic scope} of every actor. This project
targets exactly that gap; its contribution is architectural integration and a
security posture where tenant isolation is enforced by the enforcement layers
by default, not by incidental query shape.